\documentclass{article}
 
\usepackage[utf8]{inputenc}
\usepackage[T1]{fontenc}

\usepackage{times}
\usepackage{hyperref}
\usepackage[cm]{fullpage}
\usepackage{fancyhdr}
\usepackage{graphicx}
\usepackage{natbib}

\setlength{\headheight}{9.2pt}
\setlength{\headsep}{5.2pt}
\pagestyle{fancyplain}
\pagenumbering{gobble}
\lhead{\textbf{{\large Toby Dylan Hocking}  \\ \texttt{toby.hocking@mail.mcgill.ca} }}
\chead{{\LARGE \bf Projet de Recherche} }
\rhead{\includegraphics[height=0.75cm]{logo-mcgill}}

\begin{document}

\mbox{ }

Je crée de nouveaux algorithmes d'apprentissage statistique pour
résoudre des problèmes d'analyse de données dans des domaines
scientifiques tels que la médecine la biologie. Mes recherches
précédantes étaient axées sur des nouveaux modèles d'apprentissage
pour le regroupement (clustering), la régression, la détection des
ruptures, le classement (ranking) et la classification. Mes
contributions à la littérature de la statistique et l'apprentissage
automatique sont des algorithmes pour résoudre efficacement ces
problèmes d'optimisation. Je suis particulièrement intéressé à
développer des algorithmes suffisamment rapides pour l'étiquetage
interactif et l'apprentissage automatique supervisé dans les grands
ensembles de données. Ce qui suit est un résumé de mes recherches
antérieures, et mon plan de recherche pour l'avenir.

\section{Algorithmes de détection de rupture pour les données de nombre de copies d'ADN}

Pendant mon doctorat à Paris, j'ai travaillé en collaboration avec des
médecins de l'Institut Curie sur les méthodes d'apprentissage
automatique pour le diagnostic de neuroblastome, un cancer qui est
fréquent chez les enfants. Les médecins voulaient donner des
traitements ciblés basés sur le nombre de ruptures dans les données de
nombre de copies de l'ADN de la tumeur du patient. En utilisant des
algorithmes existants pour la détection des ruptures, les médecins ont
constaté des taux d'erreur trop grand (par rapport à leur
interprétation visuelle du signal et bruit dans ces données). Ma
contribution était de développer des nouveaux algorithmes basés sur
des étiquettes fournies par les médecins.

\paragraph{Nouvelle méthode d'étiquetage pour la détection supervisée des ruptures.}

Dans notre papier de \emph{BMC Bioinformatics}
\citep{HOCKING-breakpoints}, nous avons proposé une nouvelle
formulation supervisée pour le problème de détection de ruptures. Dans
notre formulation, les médecins ont d'abord étiqueté un sous-ensemble
de leurs données avec des régions positives qui contiennent des
ruptures, et des régions négatives qui ne contiennent pas de rupture.
Nous avons proposé une méthode pour l'entrainement des algorithmes
existants, en choissisant des paramètres qui minimisent le nombre de
étiquettes incorrectes. Nous avons comparé les algorithmes en
calculant le nombre de étiquettes classées correctement dans un
ensemble de test. Nous avons démontré que la meilleure méthode
existante était un modèle de détection de rupture qui maximise la
vraisemblance gaussienne.  Dans cet article, nous avons sélectionné le
nombre de changements en utilisant un fonction de pénalité univariée
qui dépend uniquement du nombre de données, et dans notre prochain
article, nous avons étudié une généralisation pour les fonctions de
pénalité multivariées.

\paragraph{Modèles de régression pour apprendre la fonction de pénalité.}

Dans un travail publié à \emph{ICML'13} \citep{HOCKING-penalties},
nous avons proposé un algorithme pour apprendre la fonction de
pénalité dans les modèles de détection de ruptures. Nous avons
démontré que l'apprentissage d'une fonction de pénalité pour
sélectionner le nombre de ruptures est un problème de régression avec
sorties censurées. Nous avons proposé d'apprendre une fonction de
pénalité linéaire, en minimisant une fonction de perte convexe qui
exploite la structure des sorties censurées. Notre contribution
principale était une nouvelle algorithme efficace pour résoudre ce
problème d'optimisation convexe. Nous avons démontré que notre
algorithme apprend une fonction de pénalité multivariée, et fournit
des taux d'erreur moins élevé que la pénalité univariée
précédente. J'ai implémenté ces algorithmes dans le paquetage
penaltyLearning dans R, que j'ai présenté lors d'un tutoriel à
\emph{useR2017} \citep{change-tutorial}. Dans des travaux plus récents
qui ont été acceptés pour publication à \emph{NIPS'17} \citep{MMIT},
nous avons proposé un algorithme pour apprendre une fonction de
penalité non-linéaire grâce à un arbre de décision.

\paragraph{Collaborations utilisant une application Web pour la
  détection supervisée de ruptures} Dans notre papier publié dans
\emph{Bioinformatics} \citep{hocking-SegAnnDB}, nous avons proposé une
application web qui facilite l'étiquetage interactif et
l'apprentissage supervisé pour l'analyse des données sur le nombre de
copies d'ADN. Du côté pratique, notre système était novateur parce
qu'il permet aux chercheurs en génomique non seulement de visualiser
leurs données, mais aussi de les étiqueter. Cette fonctionnalité est
essentielle pour les collaborations, car lorsqu'on voit une prédiction
incorrecte, on peut la corriger immédiatement en ajoutant des
étiquettes appropriées. Notre principal contribution statistique était
une algorithme de programmation dynamique efficace pour calculer les K
ruptures les plus probables pour un ensemble donné de K étiquettes
positives. Cet algorithme était essentiel pour fournir un modèle qui
correspond à n'importe quel ensemble d'étiquettes fournies par un
utilisateur. Nous avons utilisé ce système pour faciliter les
collaborations avec les scientifiques de l'Institut Curie, Paris,
France \citep{Chicard}; et Aichi Cancer Center, Nagoya, Japon
\citep{Hocking-Leukemia-2016,m14:clonal}.

\paragraph{Algorithmes optimales avec élagage des ruptures
  candidates.} Dans notre papier publié dans \emph{Statistics and
  Computing} \citep{fpop}, nous avons étudié deux techniques pour
l'élagage des ruptures candidates dans les algorithmes de détection de
ruptures optimales. Ces techniques sont très intéressantes du point de
vue computationnel : elles permettent un temps de calcul très rapide
(linéaire dans la quantité de données). Une technique est basée sur
les fonctions, et l'autre technique est basée sur les inégalités. Nous
voulions étudier quelle méthode est plus efficace. Nous avons démontré
mathématiquement que la technique fonctionnelle élague toutes les
candidates élaguées par la technique d'inégalités. Ce résultat
théorique a suggéré que l'algorithme d'élagage fonctionnel devrait
être plus rapide. Nous avons proposé le nouveau algorithme FPOP, basé
sur la technique fonctionnelle. Dans nos tests empiriques, nous avons
observé que FPOP est nettement plus rapide que plusieurs d'autres
algorithmes pour la détection de ruptures. Nous avons aussi étudié le
taux d'erreur pour la détection de ruptures, ce qui était le même que
les autres algorithmes optimales.
  
\bibliographystyle{plain}
\bibliography{TDH-refs}

\end{document}
